\documentclass[12pt]{article}
\pagestyle{empty}
\usepackage{setspace}
\usepackage{amsmath,amssymb,euscript,graphicx}
\usepackage{xurl}
\usepackage{graphicx}
\usepackage{subfigure}
\usepackage{wrapfig}
\usepackage[compact]{titlesec}
\usepackage{fancybox,color}
\usepackage[footnotesize]{caption}
\usepackage{cite}
\usepackage{enumitem}
\usepackage[normalem]{ulem}
\usepackage[top=0.9375in, right=0.9375in, left=0.9375in, bottom=0.9375in, centering]{geometry}
\usepackage{bm}

\usepackage{amsthm}


% -- Loading the code block package:
\usepackage{listings}
% -- Basic formatting
\usepackage[utf8]{inputenc}
\usepackage[english]{babel}
\usepackage{times}
\setlength{\parindent}{8pt}
\usepackage{indentfirst}
% -- Defining colors:
\usepackage[dvipsnames]{xcolor}
\definecolor{codegreen}{rgb}{0,0.6,0}
\definecolor{codegray}{rgb}{0.5,0.5,0.5}
\definecolor{codepurple}{rgb}{0.58,0,0.82}
\definecolor{backcolour}{rgb}{0.95,0.95,0.92}
% Definig a custom style:
\lstdefinestyle{mystyle}{
    backgroundcolor=\color{backcolour},   
    commentstyle=\color{codepurple},
    keywordstyle=\color{NavyBlue},
    numberstyle=\tiny\color{codegray},
    stringstyle=\color{codepurple},
    basicstyle=\ttfamily\footnotesize\bfseries,
    breakatwhitespace=false,         
    breaklines=true,                 
    captionpos=t,                    
    keepspaces=true,                 
    numbers=left,                    
    numbersep=5pt,                  
    showspaces=false,                
    showstringspaces=false,
    showtabs=false,                  
    tabsize=2
}
% -- Setting up the custom style:
\lstset{style=mystyle}


% IEEE uses Times Roman font, so we'll default to Times.
% These three commands make up the entire times.sty package.
\renewcommand{\sfdefault}{phv}
\renewcommand{\rmdefault}{ptm}
\renewcommand{\ttdefault}{pcr}

\newcommand{\squeeze}{\vspace{-0.1in}}
\newcommand{\changed}[1]{\textcolor{blue}{#1}} %generates text in blue

%%%%%%%%%%%%%%%%%%%%%%%%%%%%%%%%%%%%%%%%%%%%%%%%%%%%%%%%%%%%%%%%%%%%%%%

\begin{document}

\begin{center}
  {\Large {\bf Check-In 1}}
  
\end{center}
\vspace{5pt}
\hrule
\vspace{5pt}

Name:

\vspace{5mm}

Email:


\vspace{5pt}
\hrule
\vspace{5pt}

\begin{lstlisting}[language=Python]
def computeValue(t1, t2):
    toatl = t1 - t2
    return(total)

y1 = 5
y2 = 10
value = computeValue(y1, y2)
print("The value is ", value)

x = "fifteen"
print("The answer is: " + x)
\end{lstlisting}


\vspace{5pt}
\hrule
\vspace{5pt}

\begin{enumerate}
\item Write down a representation of the global symbol table (name, type, value) at the last line of code. The symbol table should include all identifiers defined by the code up to that point.

\vspace{45mm}
  
\item Write down what the program prints when run.

  \vspace{15mm}

\item In what symbol table do the parameters \texttt{\textcolor{red}{t1}} and \textcolor{red}{\texttt{t2}} belong?

  \vspace{15mm}

\item Would the code run if you added one or both of the following lines: \texttt{forward(100)} or \texttt{turtle.forward(100)}? Why or why not? 

  \vspace{15mm}
  
\end{enumerate}


%\Large{\bf Acknowledgments}

%\normalsize{Sources: These questions are inspired by Professor Bruce Maxwell's CS 151 and CS 152 course. The quotes and character names come from Star Trek the Next Generation.}

\end{document}
